\section*{Introduction générale}
\addcontentsline{toc}{section}{Introduction générale}

La gestion des consommations d'eau et d'électricité repose sur la collecte régulière des index des compteurs, la vérification des valeurs saisies et le calcul des montants à facturer. Une application web permet de centraliser ces opérations, de conserver les justificatifs photographiques et de réduire les erreurs de saisie.

Le projet \textbf{AquaVolt} a pour objectif de fournir une application permettant aux clients de transmettre leurs relevés mensuels et de consulter leurs factures, tandis que le fournisseur gère les abonnés, contrôle les anomalies et consulte les factures générées.

La version Alpha se concentre sur trois modules : la gestion des clients, la saisie des consommations et la gestion des factures. Les réclamations, les paiements, l'import annuel et le tableau de bord avancé sont réservés à une version ultérieure.

Le présent rapport décrit le contexte et les besoins du projet, l'architecture MVC adoptée, le modèle de données, les fonctionnalités développées, les mesures de sécurité, les tests réalisés et les perspectives d'amélioration.
